\documentclass[10pt]{article}
\usepackage[T1]{fontenc}
\usepackage[utf8]{inputenc}
\usepackage[french]{babel}
\usepackage{amsmath,amssymb,mathrsfs,mathpazo}
\usepackage{geometry}
\geometry{a4paper,margin=2cm}
\usepackage{graphicx,array,tabularx,tabularray,colortbl,multirow,multicol,booktabs,enumitem,pifont,ulem,fancybox,diagbox,cancel}
\usepackage{pstricks,pstricks-add,pst-plot,pst-tree,pst-node,pst-text,pst-func,pst-math,pst-bspline,pst-3dplot,pst-eucl,pst-solides3d}
\let\pscalculate\undefined
\usepackage{pst-bezier}
\usepackage{tikz,pgfplots,tkz-tab}
\usetikzlibrary{arrows,arrows.meta,patterns,calc,positioning,trees,decorations.markings,shapes.geometric,shapes.misc,shapes.arrows,intersections,decorations.pathmorphing,decorations.pathreplacing,backgrounds,shadows,babel}
\pgfplotsset{compat=1.18}
\usepgfplotslibrary{fillbetween}
\usepackage[np]{numprint}
\usepackage{listings,makecell,pas-tableur,twoopt}
\let\barre\undefined
\usepackage[pstricks]{tablvar}
\definecolor{darkgreen}{RGB}{0,100,0}
\providecommand{\red}{\color{red}}
\providecommand{\blue}{\color{blue}}
\providecommand{\green}{\color{green}}
\providecommand{\black}{\color{black}}
\providecommand{\white}{\color{white}}
\providecommand{\gray}{\color{gray}}
\providecommand{\cyan}{\color{cyan}}
\providecommand{\magenta}{\color{magenta}}
\providecommand{\yellow}{\color{yellow}}
\providecommand{\violet}{\color{violet}}
\providecommand{\orange}{\color{orange}}
\providecommand{\euro}{\texteuro}
\providecommand{\EUR}{\texteuro}
\providecommand{\ieme}{\textsuperscript{e}}
\providecommand{\degres}{\ensuremath{^\circ}}
\providecommand{\degree}{\ensuremath{^\circ}}
\providecommand{\vv}[1]{\overrightarrow{#1}}
\providecommand{\Coord}[1]{#1}
\providecommand{\decosix}{}
\providecommand{\starredbullet}{\ensuremath{\bullet}}
\pagestyle{empty}
\setlength{\parindent}{0pt}
\usepackage{adjustbox}\DeclareUnicodeCharacter{2236}{:}\DeclareUnicodeCharacter{2190}{\ensuremath{\leftarrow}}
\providecommand{\R}{\mathbb{R}}
\providecommand{\N}{\mathbb{N}}
\providecommand{\Z}{\mathbb{Z}}
\providecommand{\Q}{\mathbb{Q}}
\providecommand{\D}{\mathbb{D}}
\providecommand{\C}{\mathbb{C}}
\providecommand{\e}{{\mathrm{e}}}
\providecommand{\vect}[1]{\overrightarrow{#1}}
\providecommand{\vectt}[1]{\overrightarrow{\mathrm{#1}}}
\providecommand{\vvt}[1]{\overrightarrow{\mathrm{#1}}}
\providecommand{\np}[2][]{{#2~#1}}
\providecommand{\vv}[1]{{\overrightarrow{#1}}}
\renewcommand{\d}{{\mathrm{d}}}
\providecommand{\degres}{{^\circ}}
\providecommand{\degre}{{^\circ}}
\providecommand{\blue}{\color{blue}}
\providecommand{\red}{\color{red}}
\providecommand{\enskip}{\quad}
\providecommand{\llbracket}{\lbrack\!\lbrack}
\providecommand{\rrbracket}{\rbrack\!\rbrack}
\providecommand{\euro}{{\text{€}}}
\providecommand{\no}{{\text{n°}}}
\providecommand{\interval}[2]{{[#1;#2]}}
\providecommand{\intervFO}[2]{{[#1;#2[}}
\providecommand{\dotfill}{\ldots}

\providecommand{\ds}{\displaystyle}
\providecommand{\textding}[1]{#1}
\providecommand{\Oij}{\(\mathrm{O};\vec{\imath},\vec{\jmath}\)}
\providecommand{\Oijk}{\(\mathrm{O};\vec{\imath},\vec{\jmath},\vec{k}\)}
\providecommand{\Ouv}{\(\mathrm{O};\vec{u},\vec{v}\)}
\providecommand{\intervFF}[2]{\ensuremath{[#1;#2]}}
\renewcommand{\euro}{\ifmmode\text{\texteuro}\else\texteuro\fi}
\DeclareUnicodeCharacter{00B0}{\ifmmode^{\circ}\else\textdegree\fi}
\DeclareUnicodeCharacter{20AC}{\ifmmode\text{\texteuro}\else\texteuro\fi}
% Énoncés : Éducation nationale. Transcription et corrigés : APMEP, Denis Vergès et contributeurs. Découpage et classement : CoopMaths, Sébastien Lozano, Éric Elter et contributeurs. Adaptation de la présentation : Maths Explorer.
% CC BY-SA : https://coopmaths.fr/www/about/
% https://coopmaths.fr/alea/static/eam/2026/tex/eam_specifique_2026_09_metropole_3_cor.tex
\begin{document}
\noindent\textit{APMEP / CoopMaths. CC BY-SA. Présentation adaptée par Maths Explorer.}\par\medskip
\begin{enumerate}
	\item On complète le tableau par différences : lacs $= 500 - 300 = 200$ ;
	lacs satisfaisants $= 200 - 80 = 120$ ; rivières satisfaisantes $= 280 - 120 = 160$ ;
	rivières insatisfaisantes $= 300 - 160 = 140$ ; insatisfaisants $= 500 - 280 = 220$.

	\begin{center}
		\begin{tblr}{hlines, vlines,colspec={X[l,5cm]X[c,2cm]X[c,2cm]X[c,2cm]}}
			& Rivière & Lac & Total \\
			Qualité satisfaisante & \textbf{160} & \textbf{120} & 280 \\
			Qualité insatisfaisante & \textbf{140} & 80 & \textbf{220} \\
			Total & 300 & \textbf{200} & 500 \\
		\end{tblr}
	\end{center}

	\item Le choix se fait au hasard : on est en situation d'équiprobabilité.
	\begin{enumerate}
		\item $P(R) = \dfrac{300}{500} = \dfrac{3}{5}$.
		
		La probabilité que l'échantillon provienne d'une rivière est de $\frac{3}{5}$.

		\item $\overline{S}$ : \og l'échantillon choisi est classé de qualité insatisfaisante \fg.
		\[P\left(\overline{S}\right) = \dfrac{220}{500} = \dfrac{10 \times 2 \times 11}{10 \times 2 \times 25} = \dfrac{11}{25}.\]

		\item $S \cap \overline{R}$ est l'évènement \og l'échantillon est de qualité satisfaisante
		et provient d'un lac \fg.
		\[P\left(S \cap \overline{R}\right) = \dfrac{120}{500} = \dfrac{10 \times 2 \times 6}{10 \times 2 \times 25} = \dfrac{6}{25}.\]
		La probabilité que l'échantillon choisi provienne d'un lac et soit de qualité
		satisfaisante est égale à $\dfrac{6}{25}$.
	\end{enumerate}

	\item
	\begin{enumerate}
		\item $x = P_R\left(\overline{S}\right)$ est la probabilité que l'échantillon soit de qualité insatisfaisante
		sachant qu'il provient d'une rivière.

		\item Parmi les 300 échantillons provenant d'une rivière, 140 sont de qualité insatisfaisante :
		\[x = P_R\left(\overline{S}\right) = \dfrac{140}{300} = \dfrac{10 \times 2 \times 7}{10 \times 2 \times 15} = \dfrac{7}{15}.\]

		\item $P\left(\overline{R}\right) = 1 - \dfrac{3}{5} = \dfrac{2}{5}$ ; $P_R(S) = 1 - \dfrac{7}{15} = \dfrac{8}{15}$ ;
		$P_{\overline{R}}(S) = \dfrac{120}{200} = \dfrac{3}{5}$ et
		$P_{\overline{R}}\left(\overline{S}\right) = \dfrac{80}{200} = \dfrac{2}{5}$.

		\begin{center}
			\begin{tikzpicture}[x=2.8cm, y=1cm]
				
				\node (root) at (0,0) {};
				\node (R) at (1,1.3) {$R$};
				\node (Rb) at (1,-1.3) {$\overline{R}$};
				\node (S1) at (2,2.1) {$S$};
				\node (S1b) at (2,0.5) {$\overline{S}$};
				\node (S2) at (2,-0.5) {$S$};
				\node (S2b) at (2,-2.1) {$\overline{S}$};
				
				\draw (root.east) -- (R) node[pos=0.5,above] {$\frac{3}{5}$};
				\draw (root.east) -- (Rb) node[pos=0.5,below] {$\frac{2}{5}$};
				
				\draw (R.east) -- (S1) node[pos=0.5,above] {$\frac{8}{15}$};
				\draw (R.east) -- (S1b) node[pos=0.5,below] {$\frac{7}{15}$};
				\draw (Rb.east) -- (S2) node[pos=0.5,above] {$\frac{3}{5}$};
				\draw (Rb.east) -- (S2b) node[pos=0.5,below] {$\frac{2}{5}$};
			\end{tikzpicture}
		\end{center}
	\end{enumerate}

	\item Les évènements $R$ et $\overline{S}$ sont indépendants si et seulement si
	$P\left(R \cap \overline{S}\right) = P(R) \times P\left(\overline{S}\right)$.

	D'une part $P\left(R \cap \overline{S}\right) = \dfrac{140}{500} = \dfrac{7}{25} = \dfrac{35}{125}$ ;
	d'autre part $P(R) \times P\left(\overline{S}\right) = \dfrac{3}{5} \times \dfrac{11}{25} = \dfrac{33}{125}$.

	Comme $\dfrac{35}{125} \neq \dfrac{33}{125}$, l'égalité n'est pas vérifiée.

	Les évènements \og l'échantillon provient d'une rivière \fg{} et \og l'échantillon est de qualité
	insatisfaisante \fg{} ne sont \textbf{pas indépendants}.
\end{enumerate}

\end{document}
